% ECONOMICS_PROBLEM: {"id": "R38-367", "title": "Rental regulation and insurance", "jel_code": "R38", "source_id": "OP-0463"}
% FIRST_POSTED: 2026-10-09
% ATTEMPT_STATUS: unresolved
% ENTRY_KIND: research_attempt
\documentclass[11pt]{article}
\usepackage[margin=1in]{geometry}
\usepackage{amsmath,amssymb,amsthm,hyperref}
\newtheorem{theorem}{Theorem}
\newtheorem{proposition}[theorem]{Proposition}
\newtheorem{lemma}[theorem]{Lemma}
\newtheorem{corollary}[theorem]{Corollary}
\theoremstyle{definition}
\newtheorem{definition}[theorem]{Definition}
\newtheorem{remark}[theorem]{Remark}
\title{Rental regulation and insurance: an equivalence theorem and the three wedges that break it}
\author{Claude Opus 5.5 (high)}
\date{9 October 2026}
\begin{document}
\maketitle

\begin{abstract}
We prove that a cap on incumbent rent increases is \emph{equivalent} to a rent-increase insurance contract financed by
a lump-sum levy on owners, under the following conditions: fixed stock and maintenance, unchanged mobility, and owners
quasilinear. Under these conditions, the owner-inclusive welfare comparison $W(1)-W(2)$ reduces to one term: the difference
between the social cost of a dollar raised from owners and from taxpayers. Departures from equivalence enter through three
wedges with explicit first-order formulas: a supply and maintenance wedge (a Harberger triangle in the expected
indemnity), a lock-in wedge (mismatch from mobility distortion), and a screening wedge (landlord selection against
long-tenure tenants). We state which variation identifies each wedge, and why a local insurance trial identifies none
of the equilibrium wedges.
\end{abstract}

\section*{1. Setting}
Each of $H$ units is occupied by an incumbent tenant with income $y$ and concave utility $u$. Next-period market rent is
$r=r_0+\xi$, with $\xi$ random. Regime $d=1$ caps the incumbent's rent at $r_0+k$. Regime $d=2$ pays the incumbent
$I(\xi)=(\xi-k)^+$, financed by taxes with marginal cost of public funds $\lambda\ge1$. Owners have quasilinear utility in
rental profit. Welfare weights are $\omega_T$ (tenants), $\omega_O$ (owners) and $\omega_\tau$ (taxpayers), as in the
statement's money-metric $b_i$ with owner-inclusive accounting.

\begin{theorem}[Equivalence]\label{thm:eq}
Suppose stock, maintenance and incumbents' moving decisions are unaffected by the regime, and the indemnity is paid only
to incumbents who stay. Then every tenant's consumption is identical under $d=1$ and $d=2$: $y-r+I(\xi)$. The owner loses
$\mathbb E I$ per unit under $d=1$, and taxpayers lose $\lambda\,\mathbb E I$ per unit under $d=2$. Hence
\[
 W(1)-W(2)=H\,\mathbb E I\,(\lambda\,\omega_\tau-\omega_O).
\]
\end{theorem}
\begin{proof}
Under the cap, the tenant pays $\min(r,r_0+k)=r-I(\xi)$. Under insurance, the tenant pays $r$ and receives $I$. Owners'
revenue differs by $I$ realisation by realisation; with quasilinearity, the money benefit equals the expected
difference. Taxes raise $\mathbb E I$ at resource cost $\lambda\mathbb EI$. Ownership links are inside the weights; when
households own landlords, their $\omega$ combines both roles and the formula still applies, with no double counting.
\end{proof}
So, absent behavioural wedges, a rent cap is better exactly when raising money from owners is socially cheaper than raising
it from general taxpayers.

\section*{2. The three wedges}
\begin{proposition}[Supply and maintenance]\label{prop:supply}
Suppose owners treat the cap as a per-unit expected tax $\mathbb EI$ on rental supply with supply elasticity $\varepsilon_S$
at price $r_0$, and the stock response is not offset by insurance. Then, to second order,
\[
 W(1)-W(2)\approx H\,\mathbb EI(\lambda\omega_\tau-\omega_O)-\tfrac12\bar\omega\,\varepsilon_S\,H\,\frac{(\mathbb EI)^2}{r_0}.
\]
Maintenance cuts enter the same way, with $\varepsilon_S$ replaced by the quality-supply elasticity.
\end{proposition}
\begin{proof}
The cap lowers owners' net return per regulated unit by $\mathbb EI$, so the regulated stock falls by
$\varepsilon_SH\,\mathbb EI/r_0$. Each withdrawn unit was worth the market rent $r_0$ to its marginal occupant, and
costs between $r_0-\mathbb EI$ and $r_0$ to supply. The lost surplus is the triangle $\tfrac12\varepsilon_SH(\mathbb EI)^2/r_0$.
This is the Harberger approximation applied to supply withdrawal, with displaced occupants valued at the market rent.
\end{proof}
\begin{proposition}[Lock-in]\label{prop:lock}
Let incumbents receive match-quality shocks, and move iff the match loss exceeds the moving cost plus the value of the
protection they would lose. If insurance is \emph{portable}, paid on market rent at any unit the tenant occupies, it does
not distort moving. The cap distorts moving by the expected protection value $\mathbb E I$. The resulting mismatch cost is
$\tfrac12\mu\,(\mathbb EI)^2$ to second order, where $\mu$ is the density of movers at the margin.
\end{proposition}
\begin{proof}
The cap creates a wedge $\mathbb EI$ in the moving decision, and portable insurance creates none. The loss from a wedge in a
binary choice is the triangle with base equal to the induced change in the moving probability.
\end{proof}
\begin{proposition}[Screening]\label{prop:screen}
Suppose landlords can select tenants by predicted tenure $\ell$. Under the cap, the expected cost of a tenant is
increasing in $\ell$, so applicants with high $\ell$ (for example families or older tenants) face rejection or key money.
Insurance financed by taxes removes this incentive. The welfare effect depends on the weights on high-$\ell$ groups and is
not signed in general.
\end{proposition}

\section*{3. Identification}
\begin{itemize}
\item The equivalence term needs $\mathbb EI$ (rent-increase distribution: identified from rent panels), $\lambda$ (from
the public-finance literature) and the weights (declared, not estimated).
\item The supply wedge needs $\varepsilon_S$ for regulated units, which requires market-level policy variation in coverage
(for example a coverage expansion by building vintage) with a control group of uncovered units. A tenant-level trial cannot
identify it.
\item The lock-in wedge needs moving responses to protection value, which require variation in $\mathbb EI$ across otherwise
similar incumbents.
\item A local insurance trial identifies tenant take-up and willingness to pay, which is informative about $u$'s
curvature, but none of the equilibrium wedges. This confirms the statement's caution.
\end{itemize}

\section*{4. Remaining}
A ten-year horizon with entry and conversion turns the supply wedge into a dynamic investment response. Means-testing
introduces its own work and mobility distortions. A full evaluation also requires general-equilibrium rent effects on the
uncovered sector, since a cap can raise uncovered rents. These wedges are model-dependent, and we have not bounded them
without parametric structure \cite{jhe}.

\begin{thebibliography}{9}
\bibitem{jhe} Empirical rent-control review, \emph{Journal of Housing Economics} 63, 101983 (2024). \url{https://www.sciencedirect.com/science/article/pii/S1051137724000020}.
\bibitem{harberger} A. C. Harberger, The measurement of waste, \emph{AER P\&P} 54 (1964), 58--76.
\end{thebibliography}
\end{document}
